% (C) 2026 Ralf Stephan, in collaboration with Claude Code.  CC0 1.0.
\chapter{Carry graphs, walks and words}\label{ch:graphs}

This chapter fixes the language: carry graphs and their constructions, walks and infinite
paths, the words a carry graph reads, and how to count them. It ends with Kőnig's lemma, the
one place where finite branching is used in its full strength. Everything here is standard
\cite[§1.3, §2.2, §3.1]{LM95}, \cite[§1.2--1.4]{BJG09}. The only departure from the textbooks
is that the vertex set may be infinite.

\section{Carry graphs}

\begin{definition}\label{def:carry-graph}
Let \(V\) and \(A\) be sets. A \emph{carry graph}\index{carry graph} \(X\) on \(V\) over the alphabet \(A\) is a
map \(\succc_X\colon V\times A\to\{\text{finite subsets of }V\}\). We write
\(v\xrightarrow{a}w\), or \(X\colon v\xrightarrow{a}w\), when \(w\in\succc_X(v,a)\), and call it
an \emph{arc}\index{arc} with \emph{label}\index{label} \(a\). The \emph{underlying digraph}\index{underlying digraph} has an arc \(v\to w\) when
\(v\xrightarrow{a}w\) for some \(a\).
\end{definition}

When \(A\) is finite, every vertex has finitely many out-neighbours in the underlying digraph,
namely \(\operatorname{out}(v)=\bigcup_a\succc(v,a)\). In that case \(X\) is \emph{finitely
branching}, which is all that Kőnig's lemma needs. No finiteness is assumed for \(V\).
The follower automaton of Chapter~\ref{ch:sofic} is a carry graph whose vertex set is infinite
in general.

Parallel arcs with the same label are not allowed. Two arcs \(v\to w\) with distinct labels
are allowed and are counted separately. In the terms of \cite[§3.1]{LM95}, a carry graph is a
labelled graph with at most one edge per (source, label, target) triple. Every finite labelled
relation \(E\subseteq V\times A\times V\) is a carry graph, and every carry graph with finite
\(V\) and \(A\) is such a relation: its set of arcs.

\begin{example}[digit graph]\label{ex:digit}\index{digit graph}
Let \(a\ge2\) and \(M\ge1\). The \emph{digit graph} \(D_{a,M}\) has vertex set \(\Z/M\),
alphabet \(\{0,\dots,a-1\}\), and an arc \(r\xrightarrow{d}ar+d\) whenever both \(r\) and
\(ar+d\) are units. If \(x_n=\lfloor\xi a^n\rfloor\) and \(d_n\) is the \(n\)-th base-\(a\)
digit of \(\xi\), then \(x_{n+1}=ax_n+d_n\). So the residues of \(x_n\) walk along \(D_{a,M}\)
for as long as they are units, and they read the digits of \(\xi\).
\end{example}

\begin{example}[cell graph]\label{ex:cell}
Let \(p>q\ge1\), \(\nu\in\R\), \(D\ge1\), and let \(S\subset\Z\) be finite. The cell graph has
vertex set \(\{0,\dots,D-1\}\), alphabet \(S\), and an arc \(c\xrightarrow{s}c'\) when the
interval \(\bigl(p\,[c/D,(c+1)/D)-(p-q)\nu-s\bigr)/q\) meets \([c'/D,(c'+1)/D)\). By the
carry identity~\eqref{eq:carry}, the cells of the fractional parts \(y_n\) walk along it.
\end{example}

\paragraph{Constructions.} For a carry graph \(X\):
\begin{itemize}
\item the \emph{reverse} \(X^{\top}\) turns every arc around and keeps its label;
\item the \emph{restriction} \(X|_S\) to \(S\subseteq V\) keeps the arcs with both ends in
  \(S\);
\item for \(g\colon A\to B\), the \emph{relabelling} \(g_*X\) has \(v\xrightarrow{b}w\) iff
  \(v\xrightarrow{a}w\) in \(X\) for some \(a\) with \(g(a)=b\); forgetting labels is the case
  \(B=\{*\}\).
\end{itemize}
A \emph{homomorphism}\index{homomorphism} \(f\colon X\to Y\) of carry graphs over the same alphabet is a map on
vertices sending arcs to arcs with the same label:
\(v\xrightarrow{a}w\) implies \(f(v)\xrightarrow{a}f(w)\).

\begin{example}\label{ex:proj}
If \(M'\mid M\), reduction modulo \(M'\) is a homomorphism \(D_{a,M}\to D_{a,M'}\), because ring
homomorphisms send units to units and preserve \(ar+d\). This is the map behind the ladder of
Chapter~\ref{ch:reach}.
\end{example}

\section{Walks, infinite paths, words}

A \emph{walk} of length \(n\) is a sequence \(g_0\to g_1\to\dots\to g_n\) in the underlying
digraph. An \emph{infinite path} is a pair of sequences \((s,\ell)\) with
\(s_n\xrightarrow{\ell_n}s_{n+1}\) for all \(n\). We say that \emph{reading} the word
\(w=a_1\cdots a_k\in A^*\) from \(v\) can end at \(u\) if there is a walk
\(v=t_0\xrightarrow{a_1}t_1\to\dots\xrightarrow{a_k}t_k=u\). The word \(w\) is a \emph{word of
\(X\) at \(v\)} if reading it from \(v\) ends somewhere. These words form the \emph{label
language} \(\mathcal L_v(X)\).

\begin{proposition}\label{prop:words}
\begin{enumerate}[label=(\alph*)]
\item Reading \(w_1w_2\) from \(v\) ends at \(u\) iff reading \(w_1\) ends at some \(t\) from
  which reading \(w_2\) ends at \(u\).
\item \(\mathcal L_v(X)\) contains the empty word and is closed under prefixes.
\item For a finite \(S\subseteq V\), the set \(\operatorname{ends}(S,w)\) of vertices where reading
  \(w\) from \(S\) can end is finite. It is computed letter by letter by
  \(\operatorname{ends}(S,aw)=\operatorname{ends}\bigl(\bigcup_{v\in S}\succc(v,a),\,w\bigr)\),
  and \(w\in\mathcal L_v(X)\) iff \(\operatorname{ends}(\{v\},w)\neq\emptyset\).
\end{enumerate}
\end{proposition}

\begin{proof}
(a) is induction on \(w_1\). (b) follows from (a). (c) is induction on \(w\): a finite union of
finite sets is finite.
\end{proof}

Part (c) means that membership in \(\mathcal L_v(X)\) is decidable whenever equality of vertices
is, even for infinite \(V\). The sets \(\operatorname{ends}(\{v\},w)\) are the states of the
subset construction of Chapter~\ref{ch:sofic}.

\section{Counting paths and words}

Let \(A\) be finite. The number of labelled paths of length \(n\) from \(v\) is
\[
  N_v(0)=1,\qquad N_v(n+1)=\sum_{a\in A}\ \sum_{u\in\succc(v,a)}N_u(n),
\]
and \(\mathcal L_v^n\) denotes the set of words of length \(n\) at \(v\).

\begin{definition}
\(X\) is \emph{right-resolving}\index{right-resolving} if \(|\succc(v,a)|\le1\) for all \(v,a\) \cite[Def.~3.3.1]{LM95}.
\(X\) is \emph{functional}\index{functional} if at most one arc leaves each vertex.
\end{definition}

\begin{proposition}\label{prop:count}
\begin{enumerate}[label=(\alph*)]
\item \(|\mathcal L_v^n|\le N_v(n)\).
\item If \(X\) is right-resolving, then \(|\mathcal L_v^n|=N_v(n)\): a path from \(v\) is
  determined by its label \cite[Prop.~3.3.9]{LM95}.
\item If \(X\) is functional, then \(N_v(n)\le1\).
\end{enumerate}
\end{proposition}

\begin{proof}
Both sides of (a) satisfy the same recursion, with \(\le\) on the left because the words
\(a\,w'\) obtained from different successors can coincide. For (b), words with different first
letters are different. For a fixed first letter there is at most one successor, and prefixing
\(a\) is injective, so the recursion holds with equality. For (c), induct on \(n\): at most one
term of the outer sum is non-empty, and it has one summand.
\end{proof}

\section{Kőnig's lemma}

\begin{theorem}[Kőnig]\label{thm:konig}\index{K\H{o}nig's lemma}
Let \(A\) be finite and \(v\in V\). If there are walks of every length starting at \(v\), then
there is an infinite path starting at \(v\). In particular, if \(\mathcal L_v(X)\) contains
words of every length, then \(X\) has an infinite path \((s,\ell)\) with \(s_0=v\).
\end{theorem}

\begin{proof}
Call a vertex \emph{good} if walks of every length start at it. Let \(v\) be good, and suppose
no successor of \(v\) is good. For each \(w\in\operatorname{out}(v)\) choose \(N_w\) such that no
walk of length \(N_w\) starts at \(w\), and put \(N=\max_wN_w\). This maximum exists because
\(\operatorname{out}(v)\) is finite. A walk of length \(N+1\) from \(v\) passes through some
\(w\in\operatorname{out}(v)\) and continues with a walk of length \(N\ge N_w\) from \(w\). Its
prefix of length \(N_w\) is then a walk of length \(N_w\) from \(w\), a contradiction. So every
good vertex has a good successor, and choosing one repeatedly gives the infinite path. Labels
are chosen arc by arc.
\end{proof}

Mathlib proves Kőnig's lemma for strongly atomic orders and inverse systems
(\texttt{Order/KonigLemma}). For carry graphs the direct argument above is shorter than
either reduction. A sharper version, with an injective path, is Theorem~\ref{thm:geodesic-konig}
in the next chapter.

\section*{Notes}
The definitions follow \cite[§2.2, §3.1]{LM95} and \cite[§1.2]{BJG09}. The only difference is
that vertex sets may be infinite, with finitely many arcs per label. Right-resolving
presentations are the subject of \cite[§3.3]{LM95}. The digit graph is the coprimality subshift
of \cite{FS67,Dub09}, and the cell graph goes back to \cite{FLP95}.

\section*{Lean dictionary}
\begin{center}\small
\begin{tabular}{@{}>{\raggedright}p{0.3\textwidth}>{\raggedright\arraybackslash}p{0.64\textwidth}@{}}
\toprule
Text & Lean (\texttt{CarryGraph/Basic.lean}, \texttt{Path.lean})\\
\midrule
Definition~\ref{def:carry-graph} & \lean{CarryGraph}, \lean{Adj}, \lean{under}, \lean{out}\\
finite relations, arc sets & \lean{ofRel}, \lean{ofFinset}, \lean{arcs}, \lean{ofFinset\_arcs}\\
constructions & \lean{reverse}, \lean{restrict}, \lean{relabel}, \lean{forget}, \lean{Hom}\\
Example~\ref{ex:digit}, \ref{ex:proj} & \lean{CarryBridge.digitGraph}, \lean{CarryBridge.projHom}\\
walks, infinite paths & \lean{IsWalk}, \lean{IsIPath}\\
Proposition~\ref{prop:words} & \lean{Reads}, \lean{IsWord}, \lean{reads\_append},
  \lean{IsWord.of\_append}, \lean{ends}, \lean{isWord\_iff\_ends\_nonempty}\\
Proposition~\ref{prop:count} & \lean{npaths}, \lean{words}, \lean{card\_words\_le\_npaths},
  \lean{card\_words\_eq\_npaths}, \lean{Functional.npaths\_le\_one}\\
Theorem~\ref{thm:konig} & \lean{exists\_walk\_of\_forall\_hasWalk},
  \lean{exists\_iPath\_of\_forall\_isWord}\\
\bottomrule
\end{tabular}
\end{center}
